% Trevor Chen -- CV. Compile with: xelatex Trevor_Chen_CV.tex
\documentclass[letterpaper,10pt]{article}

\usepackage{fontspec}            % Latin Modern by default
\usepackage[left=0.7in,right=0.7in,top=0.55in,bottom=0.5in]{geometry}
\usepackage{titlesec}
\usepackage{enumitem}
\usepackage[hidelinks]{hyperref}
\usepackage{amsmath}
\usepackage{ragged2e}

\pagestyle{empty}
\linespread{1.0302}
\setlength{\tabcolsep}{0in}
\setlength{\parindent}{0pt}
\setlength{\parskip}{3pt}
\setlength{\RaggedRightParindent}{0pt}
\setlength{\JustifyingParindent}{0pt}
\raggedright

% Section headings: bold, large, ruled
\newif\ifafterhead
\titleformat{\section}{\large\bfseries}{}{0em}{}[\titlerule\global\afterheadtrue]
\titlespacing*{\section}{0pt}{8pt}{1.8pt}

% ---------- macros ----------
% Entry heading: bold name / date, italic role / place
\newcommand{\cvEntry}[4]{%
  \par\ifafterhead\vspace{5.2pt}\else\vspace{5pt}\fi\global\afterheadfalse\hspace*{0.15in}%
  \begin{tabular*}{0.97\textwidth}[t]{l@{\extracolsep{\fill}}r}
    \textbf{#1} & #2 \\
    \textit{\small #3} & \textit{\small #4} \\
  \end{tabular*}\par\nopagebreak\vspace{-0.2pt}\nopagebreak}
% Entry heading with a plain second line (service entries)
\newcommand{\cvPlainEntry}[3]{%
  \par\ifafterhead\vspace{5.2pt}\else\vspace{5pt}\fi\global\afterheadfalse\hspace*{0.15in}%
  \begin{tabular*}{0.97\textwidth}[t]{l@{\extracolsep{\fill}}r}
    \textbf{#1} & #2 \\
    \makebox[0pt][l]{\small #3} & \\
  \end{tabular*}\par\vspace{3pt}}
% Project heading: one line
\newcommand{\cvProject}[2]{%
  \par\ifafterhead\vspace{5.2pt}\else\vspace{5pt}\fi\global\afterheadfalse\hspace*{0.15in}%
  \begin{tabular*}{0.97\textwidth}{l@{\extracolsep{\fill}}r}
    #1 & #2 \\
  \end{tabular*}\par\nopagebreak\vspace{0pt}\nopagebreak}
% Indented body paragraphs under an entry
\makeatletter\newenvironment{cvBody}{\@beginparpenalty=10000\begin{list}{}{\leftmargin=0.3in\rightmargin=0pt\topsep=0pt\parsep=3pt\itemsep=0pt\partopsep=0pt}\item\relax}{\end{list}\vspace{6pt}}\makeatother
% Dated line (honors, posters): wrapped text with year on the right
\newcommand{\cvDated}[2]{%
  \par\vspace{0pt}\global\afterheadfalse\hspace*{0.15in}%
  \begin{tabular*}{0.97\textwidth}[t]{@{}p{0.84\textwidth}@{\extracolsep{\fill}}r@{}}
    \justifying #1 & #2 \\
  \end{tabular*}\par}
% Indented item without a date (publications, talks)
\newcommand{\cvItem}[1]{\par\vspace{0pt}{\leftskip=0.15in #1\par}}
% Small italic sub-label
\newcommand{\cvLabel}[1]{\par\ifafterhead\vspace{0.2pt}\else\vspace{2pt}\fi\global\afterheadfalse{\small\textit{#1}}\par\vspace{-1.6pt}}

\begin{document}

% ---------- header ----------
\begin{center}
  {\Huge\textbf{Trevor Chen}} \\ \vspace{5.6pt}
  {\small California Institute of Technology, Pasadena, CA} \\ \vspace{3pt}
  {\small +1-858-366-3389 $|$ \href{mailto:trevorbchen@gmail.com}{trevorbchen@gmail.com} $|$
  \href{https://trevorbchen.github.io/Website/}{Website} $|$
  \href{https://github.com/trevorbchen}{GitHub} $|$
  \href{https://scholar.google.com/citations?user=N_DNZMUAAAAJ}{Scholar} $|$
  \href{https://www.linkedin.com/in/trevorbchen/}{LinkedIn}}
\end{center}
\vspace{-9.9pt}

% ---------- research interests ----------
\section{Research Interests}
How to explore and exploit efficiently when experiments are expensive, how to make reinforcement learning work in AI for science, and how to interpret why models fail. My current work is on diffusion models, closed-loop discovery with wet-lab feedback, and sparse-autoencoder interpretability.

% ---------- education ----------
\section{Education}
\cvEntry{California Institute of Technology}{Pasadena, CA}
  {B.S. in Applied and Computational Mathematics $|$ GPA: 4.0/4.0}{Expected Jun. 2029}
\begin{cvBody}
Coursework: Agentic Systems, Diffusion Models, Generative Models, Deep Learning, ML \& Data Mining, GPU Programming (CUDA), Operating Systems, Statistical Inference, Probability Models, Applied Linear Algebra.
\end{cvBody}

\cvEntry{Westview High School}{San Diego, CA}
  {Graduated; dual enrollment at Palomar College (Multivariable Calculus, Linear Algebra, Differential Equations)}{2025}
\begin{cvBody}
Online coursework: Harvard CS50, Princeton Algorithms I and II, and the Stanford Machine Learning Specialization.
\end{cvBody}

% ---------- honors ----------
\section{Honors and Awards}
\cvDated{Regeneron Science Talent Search Scholar}{2025}
\cvDated{First Award, Greater San Diego Science \& Engineering Fair (Senior Division, Physics), for \textit{Perfecting Infrared Gratings for Polarized and Non-Polarized Applications}; DRS Special Award, sponsored by DRS Daylight Solutions}{2025}
\cvDated{First Place, Greater San Diego Science \& Engineering Fair (Senior Division, Physics); DRS Special Award (1 of about 200 students); 4th Place, Physics \& Astronomy, California Science and Engineering Fair}{2024}
\cvDated{USACO Platinum; AIME Qualifier; National Merit Scholar; International Logic Olympiad, Silver Scholar (top 10\%)}{2024}
\cvDated{Outstanding Presentation Award, 48th Dayton-Cincinnati Aerospace Sciences Symposium (2 of about 150 attendees; the only high school student)}{2023}
\cvDated{First Place, Greater San Diego Science \& Engineering Fair (Senior Division, Biology)}{2023}

% ---------- funding ----------
\section{Funding and Fellowships}
\cvDated{Summer Undergraduate Research Fellowship (SURF), Caltech}{2026}
\cvDated{SPIE Photonics and Education Scholarship, \$3,000 (the only high school student to receive it)}{2023}

% ---------- research experience ----------
\section{Research Experience}
\cvEntry{ETH Zurich AI Center, Krause Group \normalfont{\small (joint with Arnold Lab, Caltech)}}{Oct. 2026 -- Present}
  {Researcher, with Prof. Andreas Krause and Prof. Frances Arnold}{Remote}
\begin{cvBody}
I am building a closed-loop discovery system for enzyme and protein design. A language model proposes mechanistic hypotheses, each hypothesis is distilled into a flow model that generates candidate designs, and wet-lab results confirm or refute it. Because campaigns are limited to about 10$^2$ wet-lab measurements, I am developing signed distillation, which folds confirmed hypotheses into the generator and suppresses refuted ones while staying anchored to the previous model. The system is evaluated against active search, reward-based fine-tuning and agentic LLM loops, with wet-lab validation at the Arnold Lab.
\end{cvBody}

\cvEntry{\href{https://www.yisongyue.com/research}{Yue Lab, Caltech}}{Jan. 2026 -- Present}
  {Undergraduate Researcher, advised by Prof. Yisong Yue}{Pasadena, CA}
\begin{cvBody}
I led an 11-author study with Genesis Molecular AI on online adaptation of discrete diffusion models for molecular optimization, combining RL fine-tuning, Thompson sampling, CVaR shaping and density entropy regularization. Across six molecule tasks and three protein-fitness tasks with five seeds each, the method matches the best inference-time search baseline with 58 times fewer oracle calls on molecules and 16 times fewer on proteins (NeurIPS 2026).

I also studied why extra latent dimensions delay memorization in diffusion models. With the intrinsic dimension fixed at 5, widening the latent space from 5 to 40 dimensions reduces memorization from 30\% to 0.1\%, and the onset of score overfitting is delayed by more than 25 times out to width 240. A spectral analysis of the score network's feature-correlation matrix explains the effect, and on real image latents memorization begins 3.6 times later on CelebA and 1.4 times later on CIFAR-10 (ICML 2026 HiLD; ICLR 2027, under review). Both projects ran on multi-node GPU clusters with Slurm and fault-tolerant scheduling.
\end{cvBody}

\cvEntry{Caltech SURF Fellow, Perona Vision Lab}{Jun. 2026 -- Aug. 2026}
  {Summer Undergraduate Research Fellow, advised by Prof. Pietro Perona}{Pasadena, CA}
\begin{cvBody}
I found a family of training objectives for sparse autoencoders (SAEs) that choose which features the SAE learns, so that the dictionary can be picked on purpose instead of being left to the random seed. Linear probes, planted-signal recovery and hard-image overlap show that dictionaries from different objectives differ measurably, and combining four objectives in one model costs about 1.3 points of reconstruction (CVPR 2027, under review). Separately, I trained a diffusion transformer from scratch and 36 SAEs with more than 220,000 features on every layer's residual stream and attention and MLP outputs, then ablated features across denoising steps to trace how they drive generation.
\end{cvBody}

\cvEntry{\href{https://www.sdsc.edu/}{UCSD San Diego Supercomputer Center}}{Oct. 2020 -- May 2025}
  {Computational Biology Research Intern}{San Diego, CA}
\begin{cvBody}
I built machine-learning models for drug discovery. An ensemble model with 90\% cross-validation accuracy found aldose reductase inhibitors for diabetic cataracts, and docking validated the predictions and surfaced three novel compounds predicted to outperform known treatments. With a Stanford undergraduate I built a virtual-screening model for HIV-1 protease inhibitors, also at 90\% cross-validation accuracy, and used docking simulations to correct misclassified examples, which raised accuracy by 20\%. Across these projects I screened more than 50,000 compounds on HPC clusters.

I also carried out a meta-analysis of how anelloviruses affect known cancer pathways, which found oncogenic interactions linked to cervical and liver cancers, and I mentored three interns in machine learning and in running biological simulations, using docking to validate their models.
\end{cvBody}

\cvEntry{\href{https://spectralenergies.com/}{Spectral Energies}}{Feb. 2022 -- Jul. 2025}
  {R\&D Intern: gratings (2023--2025) and borescope optics (2022)}{Remote}
\begin{cvBody}
I designed broadband mid-IR diffraction gratings for ultrafast lasers with over 90\% efficiency from 3.6 to 6.6\,\textmu m using global optimization; prototypes measured 92\% against a theoretical 94.2\%. With Penn State Nanofab I co-developed the fabrication process, combining electron-beam lithography and anisotropic wet etching to produce smooth blazed grooves, which improved efficiency by 30\%, and I tuned Au--Al$_2$O$_3$ coatings for laser damage threshold. In simulation I found double-blazing and sawtooth groove features that widened bandwidth by 50\% and raised efficiency by 80\%, respectively, and the simulation pipelines I built cut prototype cycles by 70\% and manufacturing error by 80\%.

I delivered 10 gratings to five or more national and university labs, including the Air Force Directed Energy Directorate, the Air Force Research Laboratory, the Naval Research Laboratory and the University of Ottawa, which supported more than \$1.3 million in R\&D contracts, and I wrote the internal guidebook that underpins commercialization of the grating technology. Earlier, as primary lens designer for the LWIR borescope, I developed optimization algorithms for a lens system with ten times the spatial resolution of competing borescopes, and per-pixel blackbody calibration gave temperature readings accurate to 98\%.
\end{cvBody}

% ---------- industry ----------
\section{Industry Experience}
\cvEntry{\href{https://ascade.co/}{Ascade} (Accepted to Y Combinator F25)}{Jun. 2025 -- May 2026}
  {Founder \& CTO}{San Diego, CA}
\begin{cvBody}
I built custom tooling and internal verifiable-reward environments (RLVR) for long-horizon tasks and trained an 8B model on them with GRPO; it beats Claude Opus 4.6 and other frontier models on public formal-verification benchmarks.

I also built an autoformalization pipeline that translates English requirements into formal specifications, which Z3 checks in a counterexample-guided refinement loop. We used it to ship a formally verified payment processor that reached \$25K MRR and to win a DARPA ERIS contract and other federal work on verified access control for AI agents.
\end{cvBody}

\cvEntry{\href{https://www.daylightsolutions.com/}{Leonardo DRS Daylight Solutions}}{Jun. 2025 -- Aug. 2025}
  {Machine Learning Engineer}{San Diego, CA}
\begin{cvBody}
I built a filtering model with a custom cost function that cut water-vapor noise in laser spectroscopy data by more than 98\%, and I replaced the cost function with a piecewise approximation that made the software 50\% faster than the legacy version. I shipped the model in Daylight's MIRcat laser software, where it improved signal-to-noise by 99\% in humid conditions, and I built custom software for Daylight's laser research team, including user interfaces and real-time sensor analysis.
\end{cvBody}

% ---------- publications ----------
\section{Publications}
{\small $^*$\textit{Equal contribution (shared first author)}}\par\vspace{-4.8pt}
\cvLabel{Accepted and published}
\cvItem{\textbf{T. Chen}$^*$, A. Dai$^*$, J. Yang, R. De Santi, D. Khalil, W. Chu, N. Gruver, P. Murugan, A. Goldberg, M. Al-Shedivat, Y. Yue. \textit{\href{https://arxiv.org/abs/2607.02834}{On the Design Space of Discrete Diffusion Online Adaptation for Molecular Optimization}}. \textbf{NeurIPS 2026}.}
\cvItem{\textbf{T. Chen}, R. Shahbaba, K. Peng, A. Garg. \textit{\href{https://icml.cc/virtual/2026/72322}{How Excess Latent Dimensionality Delays Memorization in Diffusion Models}}. \textbf{ICML 2026 Workshop on High-dimensional Learning Dynamics (HiLD)}.}
\cvItem{J. Y. Tang$^*$, \textbf{T. Chen}$^*$, V. L. Kouznetsova, I. F. Tsigelny. \textit{\href{https://doi.org/10.1093/infdis/jiae626}{Anelloviruses and Cancer}}. \textbf{The Journal of Infectious Diseases} 231(2), 2025.}
\cvItem{\textbf{T. Chen}, B. Liu, J. A. McCoy, et al. \textit{\href{https://opg.optica.org/ome/fulltext.cfm?uri=ome-14-5-1336&id=549158}{Development of Broadband High-Efficiency Mid-IR Gratings for High-Energy Ultrafast Lasers}}. \textbf{Optical Materials Express} 14(5), 2024.}
\cvItem{\textbf{T. Chen}, P. S. Hsu. \textit{\href{https://www.spiedigitallibrary.org/conference-proceedings-of-spie/13116/1311609/Enabling-ultra-broadband-high-efficiency-gratings-for-single-cycle-high/10.1117/12.3032561.short}{Enabling Ultra-Broadband High-Efficiency Gratings for Single-Cycle Infrared Lasers}}. \textbf{SPIE Nanoengineering}, 2024.}
\cvItem{P. S. Hsu$^*$, \textbf{T. Chen}$^*$, A. Roy, et al. \textit{\href{https://iopscience.iop.org/article/10.1088/1361-6501/acc048}{A Compact LWIR Borescope Sensor for 2D Engine Component Surface Temperature Measurement}}. \textbf{Measurement Science and Technology} 34(6), 2023.}
\cvItem{A. Roy, P. S. Hsu, A. Zhang, \textbf{T. Chen}, et al. \textit{\href{https://doi.org/10.2514/6.2023-3580}{Development of Compact LWIR Borescope Sensor for Non-contact 2D Engine Component Surface Temperature Measurement}}. \textbf{AIAA AVIATION Forum}, 2023.}
\cvItem{\textbf{T. Chen}, A. Zhang, A. Roy, P. S. Hsu, O. Bibik, S. Adhikari. \textit{\href{https://doi.org/10.1117/12.2672622}{2D High Temperature Sensor Development Using LWIR Borescope}}. \textbf{SPIE Infrared Sensors, Devices, and Applications XIII}, 2023.}
\cvItem{\textbf{T. Chen}, R. Chen, A. You, V. L. Kouznetsova, I. F. Tsigelny. \textit{\href{https://doi.org/10.1016/j.aopr.2023.09.002}{Search of Inhibitors of Aldose Reductase for Treatment of Diabetic Cataracts Using Machine Learning}}. \textbf{Advances in Ophthalmology Practice and Research} 3(4), 2023.}
\cvItem{A. K. Gao, \textbf{T. Chen}, V. L. Kouznetsova, I. F. Tsigelny. \textit{\href{https://doi.org/10.1016/j.aichem.2023.100014}{Machine-Learning-Based Virtual Screening and Ligand Docking Identify Potent HIV-1 Protease Inhibitors}}. \textbf{Artificial Intelligence Chemistry} 1(2), 2023.}
\cvItem{\textbf{T. Chen}, P. S. Hsu. \textit{\href{https://opg.optica.org/abstract.cfm?uri=LACSEA-2022-LM1B.5}{Compact Borescope Design for 2D Long-Wavelength IR Thermal Imaging Sensing}}. \textbf{Optica Optical Sensors and Sensing Congress (LACSEA)}, 2022.}
\cvLabel{Under review}
\cvItem{\textbf{T. Chen}, R. Shahbaba. \textit{How Excess Latent Dimensionality Delays Memorization in Diffusion Models}. \textbf{ICLR 2027}.}
\cvItem{\textbf{T. Chen}, I. Demler, P. Perona. \textit{Objective Geometry Steers Sparse Autoencoder Vocabularies}. \textbf{CVPR 2027}.}

% ---------- presentations ----------
\section{Presentations and Invited Talks}
\cvLabel{Posters}
\cvDated{ICML 2026 Workshop on High-dimensional Learning Dynamics (HiLD): \textit{How Excess Latent Dimensionality Delays Memorization in Diffusion Models}}{2026}
\cvDated{SPIE Nanoengineering: \textit{Enabling Ultra-Broadband High-Efficiency Gratings for Single-Cycle Infrared Lasers}}{2024}
\cvDated{48th Dayton-Cincinnati Aerospace Sciences Symposium (Outstanding Presentation Award)}{2023}
\cvDated{AIAA AVIATION Forum: \textit{Development of Compact LWIR Borescope Sensor for Non-contact 2D Engine Component Surface Temperature Measurement}}{2023}
\cvDated{SPIE Infrared Sensors, Devices, and Applications XIII: \textit{2D High Temperature Sensor Development Using LWIR Borescope}}{2023}
\cvDated{Optica Optical Sensors and Sensing Congress: \textit{Compact Borescope Design for 2D Long-Wavelength IR Thermal Imaging Sensing}}{2022}
\cvLabel{Invited talks}
\cvDated{Leonardo DRS Daylight Solutions Executive Team: I was invited by the former VP of Engineering and Advanced Technology after the DRS Special Award and presented the grating research and its potential impact.}{}
\cvDated{Regional iGEM symposium: I was invited to speak on ethics to a large panel of mostly student peers.}{}

% ---------- projects ----------
\section{Projects}
\cvProject{\textbf{Meta-Interpretability Benchmark} $|$ \small{Independent Research}}{Nov. 2026 -- Present}
\begin{cvBody}
I am building procedurally generated rule-based games where small transformers learn a known internal state, with planted deviations (shortcuts, triggers, backdoors) and clean controls. An observer model has to recover what a network learned and where it lives, graded by causal interventions against black-box baselines.
\end{cvBody}

\cvProject{\textbf{Breaking Pangram's AI-Image Detector} $|$ \small{Independent Red-Teaming}}{Aug. 2026}
\begin{cvBody}
I attacked a commercial AI-image detector in both directions, using only the free plan because API access was never granted. Resizing a real photo from 128 by 128 to 512 by 512 with an ordinary resampler was enough for the detector to call it AI, and so was re-encoding it through a variational autoencoder. Matching the frequency statistics of real images did not fool it, which suggests it relies on local patterns and not the global spectrum.

For the other direction I trained an autoencoder on CelebA with an adversarial discriminator term. Passing an AI-generated image through it produces a near-identical copy that the detector calls real, and interpolating between the original and the re-encoded image lets the verdict be set to any confidence in between. A single image can be optimized in minutes on a laptop, and a model for a whole domain trains in about a day. I disclosed the findings to Pangram's co-founder, responded to his question about how a taxonomy should separate AI-generated, AI-edited and merely resized images, and withheld the attack code because both directions enable deniability.
\end{cvBody}

\cvProject{\textbf{LLM Inference Engine from Scratch} $|$ \small{PyTorch, Triton}}{Oct. 2026}
\begin{cvBody}
I built a vLLM-style serving engine from scratch. The KV cache is cut into blocks of 16 tokens with a per-request block table, a continuous-batching scheduler prefills waiting requests when the cache has room and otherwise decodes every running sequence, and when memory runs out it preempts the most recently admitted request and recomputes it later. Its outputs match Hugging Face token for token on a small random Llama in full precision, including under memory pressure that forces eight preemptions, and 6 of 6 sampled requests match in fp16 on the real model.

I benchmarked it against Hugging Face static batching of 8, 16 and 32 on TinyLlama-1.1B with 64 synthetic requests. Static batching wastes 45 to 54 percent of its decode work on sequences that have already finished. On one A100 the engine serves 1.26 times the throughput of the best static batch and cuts median latency under load from 3.6--9.2 seconds to 0.7 seconds; on an RTX 4060 the figures are 1.44 times and 1.2 seconds. I also wrote fused Triton kernels (FlashAttention-2 prefill, paged-attention decode, add and RMSNorm, RoPE with cache write, SwiGLU) with a float32-reference test harness, and I published the recorded runs as an interactive replay on my website.
\end{cvBody}

\cvProject{\textbf{Real-Time CUDA Black Hole Renderer} $|$ \small{CS 179 GPU Programming}}{Jun. 2026}
\begin{cvBody}
With a teammate I built a CUDA/C++ ray tracer that bends light through Schwarzschild spacetime per pixel and renders gravitational lensing, Einstein rings and a Doppler-shifted accretion disk in real time. Because the spacetime is spherically symmetric, each photon stays in one plane, so every ray integrates only its radius, angle and radial velocity with an adaptive RK4 step that shrinks near the photon sphere. My teammate wrote the geodesic integrator, and I wrote the camera, the shading and the CUDA rendering and post-processing kernels (bloom and HDR tone mapping), with one thread per pixel, plus a browser viewer with live controls and FPS profiling.

We validated the renderer against general relativity with a 12-test physics suite and CPU--GPU parity tests. Weak-field light deflection is within 0.44\% of theory, photons at the critical impact parameter of $3\sqrt{3}M$ are captured or escape correctly, energy and angular-momentum drift is about $10^{-6}$, ray paths retrace within 0.16\% under time reversal, and spherical symmetry holds to $1.8\times10^{-7}$.
\end{cvBody}

\cvProject{\textbf{Trading on Binary Prediction Markets} $|$ \small{Independent Project}}{May 2025 -- Aug. 2025}
\begin{cvBody}
In summer 2025 I built a trading system for binary prediction markets, which were fragmented and thin and often priced related contracts inconsistently. It combined an RL inventory-control policy that cut execution slippage by 30\%, sentiment signals from text, arbitrage and hedging across related markets, and a secure execution API with local signing that made the trading loop six times faster.

Over a three-month stretch it returned about 8\% a week, which was a small experiment and not a result to extrapolate. The simple strategies did most of the work: buying the obviously mispriced side of inconsistent contracts and hedging captured most of the edge, and the RL and sentiment components added only a thin sliver of extra return for their complexity.
\end{cvBody}

% ---------- teaching ----------
\section{Teaching, Mentoring and Service}
\cvEntry{Project Binary (501(c)(3) nonprofit)}{Apr. 2023 -- Jun. 2025}
  {Founder and President}{}
\begin{cvBody}
I taught programming through a free curriculum that grew from a few students on Zoom to over 100 low-income students and a team of 30 volunteers, and six other communities now run the program. Through classes, grants and fundraisers we raised over \$10,000 in 18 months, and a partnership with Computer2Kids on the Technological Divide Scholarship let us donate \$16,000 of laptops to 50 families; the project also won Google's Non-Profit Ad Grant, worth \$10,000 per month in search ads. I also started and mentored robotics teams at a Title I middle school, and students went on to regional robotics championships and USACO competitions.
\end{cvBody}

\cvPlainEntry{Mentor, UCSD San Diego Supercomputer Center}{}
  {I mentored three interns in machine learning and biological simulation.}
\cvPlainEntry{Executive Tutoring Lead, Westview Peer Counseling}{Dec. 2021 -- Jan. 2025}
  {I organized weekly tutoring sessions at the library and received the PTSA Student Staff Recognition Award.}
\cvPlainEntry{Teaching Assistant, Camp Invention}{Jul. 2021}
  {I designed hands-on science lessons for children aged 6 to 10.}

% ---------- skills ----------
\section{Skills}
\cvItem{\textbf{Languages:} Python, C++, C, CUDA, Triton, SQL, Dafny, Lean 4; fluent in Mandarin\\
\textbf{Technical:} PyTorch, Transformers, Diffusion Models, Reinforcement Learning (GRPO), Sparse Autoencoders, Distributed Training, Slurm, Linux, Git, Formal Verification\\
\textbf{Code:} \href{https://github.com/trevorbchen/DDAMO}{github.com/trevorbchen/DDAMO} (NeurIPS 2026), \href{https://github.com/trevorbchen/CS_179_final}{CS\_179\_final} (black hole renderer), \href{https://github.com/trevorbchen/sae_families}{sae\_families} (SAE objectives)}
\vspace{10pt}

% ---------- references ----------
\section{References}
References are available on request.

\end{document}
